\IfFileExists{stacks-project.cls}{%
\documentclass{stacks-project}
}{%
\documentclass{amsart}
}

% For dealing with references we use the comment environment
\usepackage{verbatim}
\newenvironment{reference}{\comment}{\endcomment}
%\newenvironment{reference}{}{}
\newenvironment{slogan}{\comment}{\endcomment}
\newenvironment{history}{\comment}{\endcomment}

% For commutative diagrams we use Xy-pic
\usepackage[all]{xy}

% We use 2cell for 2-commutative diagrams.
\xyoption{2cell}
\UseAllTwocells

% We use multicol for the list of chapters between chapters
\usepackage{multicol}

% This is generally recommended for better output
\usepackage{lmodern}
\usepackage[T1]{fontenc}

% For cross-file-references
\usepackage{xr-hyper}

% Package for hypertext links:
\usepackage{hyperref}

% For any local file, say "hello.tex" you want to link to please
% use \externaldocument[hello-]{hello}
\externaldocument[introduction-]{introduction}
\externaldocument[conventions-]{conventions}
\externaldocument[sets-]{sets}
\externaldocument[categories-]{categories}
\externaldocument[topology-]{topology}
\externaldocument[sheaves-]{sheaves}
\externaldocument[sites-]{sites}
\externaldocument[stacks-]{stacks}
\externaldocument[fields-]{fields}
\externaldocument[algebra-]{algebra}
\externaldocument[brauer-]{brauer}
\externaldocument[homology-]{homology}
\externaldocument[derived-]{derived}
\externaldocument[simplicial-]{simplicial}
\externaldocument[more-algebra-]{more-algebra}
\externaldocument[smoothing-]{smoothing}
\externaldocument[modules-]{modules}
\externaldocument[sites-modules-]{sites-modules}
\externaldocument[injectives-]{injectives}
\externaldocument[cohomology-]{cohomology}
\externaldocument[sites-cohomology-]{sites-cohomology}
\externaldocument[dga-]{dga}
\externaldocument[dpa-]{dpa}
\externaldocument[sdga-]{sdga}
\externaldocument[hypercovering-]{hypercovering}
\externaldocument[schemes-]{schemes}
\externaldocument[constructions-]{constructions}
\externaldocument[properties-]{properties}
\externaldocument[morphisms-]{morphisms}
\externaldocument[coherent-]{coherent}
\externaldocument[divisors-]{divisors}
\externaldocument[limits-]{limits}
\externaldocument[varieties-]{varieties}
\externaldocument[topologies-]{topologies}
\externaldocument[descent-]{descent}
\externaldocument[perfect-]{perfect}
\externaldocument[more-morphisms-]{more-morphisms}
\externaldocument[flat-]{flat}
\externaldocument[groupoids-]{groupoids}
\externaldocument[more-groupoids-]{more-groupoids}
\externaldocument[etale-]{etale}
\externaldocument[chow-]{chow}
\externaldocument[intersection-]{intersection}
\externaldocument[pic-]{pic}
\externaldocument[weil-]{weil}
\externaldocument[adequate-]{adequate}
\externaldocument[dualizing-]{dualizing}
\externaldocument[duality-]{duality}
\externaldocument[discriminant-]{discriminant}
\externaldocument[derham-]{derham}
\externaldocument[local-cohomology-]{local-cohomology}
\externaldocument[algebraization-]{algebraization}
\externaldocument[curves-]{curves}
\externaldocument[resolve-]{resolve}
\externaldocument[models-]{models}
\externaldocument[functors-]{functors}
\externaldocument[equiv-]{equiv}
\externaldocument[pione-]{pione}
\externaldocument[etale-cohomology-]{etale-cohomology}
\externaldocument[proetale-]{proetale}
\externaldocument[relative-cycles-]{relative-cycles}
\externaldocument[more-etale-]{more-etale}
\externaldocument[trace-]{trace}
\externaldocument[crystalline-]{crystalline}
\externaldocument[spaces-]{spaces}
\externaldocument[spaces-properties-]{spaces-properties}
\externaldocument[spaces-morphisms-]{spaces-morphisms}
\externaldocument[decent-spaces-]{decent-spaces}
\externaldocument[spaces-cohomology-]{spaces-cohomology}
\externaldocument[spaces-limits-]{spaces-limits}
\externaldocument[spaces-divisors-]{spaces-divisors}
\externaldocument[spaces-over-fields-]{spaces-over-fields}
\externaldocument[spaces-topologies-]{spaces-topologies}
\externaldocument[spaces-descent-]{spaces-descent}
\externaldocument[spaces-perfect-]{spaces-perfect}
\externaldocument[spaces-more-morphisms-]{spaces-more-morphisms}
\externaldocument[spaces-flat-]{spaces-flat}
\externaldocument[spaces-groupoids-]{spaces-groupoids}
\externaldocument[spaces-more-groupoids-]{spaces-more-groupoids}
\externaldocument[bootstrap-]{bootstrap}
\externaldocument[spaces-pushouts-]{spaces-pushouts}
\externaldocument[spaces-chow-]{spaces-chow}
\externaldocument[groupoids-quotients-]{groupoids-quotients}
\externaldocument[spaces-more-cohomology-]{spaces-more-cohomology}
\externaldocument[spaces-simplicial-]{spaces-simplicial}
\externaldocument[spaces-duality-]{spaces-duality}
\externaldocument[formal-spaces-]{formal-spaces}
\externaldocument[restricted-]{restricted}
\externaldocument[spaces-resolve-]{spaces-resolve}
\externaldocument[formal-defos-]{formal-defos}
\externaldocument[defos-]{defos}
\externaldocument[cotangent-]{cotangent}
\externaldocument[examples-defos-]{examples-defos}
\externaldocument[algebraic-]{algebraic}
\externaldocument[examples-stacks-]{examples-stacks}
\externaldocument[stacks-sheaves-]{stacks-sheaves}
\externaldocument[criteria-]{criteria}
\externaldocument[artin-]{artin}
\externaldocument[quot-]{quot}
\externaldocument[stacks-properties-]{stacks-properties}
\externaldocument[stacks-morphisms-]{stacks-morphisms}
\externaldocument[stacks-limits-]{stacks-limits}
\externaldocument[stacks-cohomology-]{stacks-cohomology}
\externaldocument[stacks-perfect-]{stacks-perfect}
\externaldocument[stacks-introduction-]{stacks-introduction}
\externaldocument[stacks-more-morphisms-]{stacks-more-morphisms}
\externaldocument[stacks-geometry-]{stacks-geometry}
\externaldocument[moduli-]{moduli}
\externaldocument[moduli-curves-]{moduli-curves}
\externaldocument[examples-]{examples}
\externaldocument[exercises-]{exercises}
\externaldocument[guide-]{guide}
\externaldocument[desirables-]{desirables}
\externaldocument[coding-]{coding}
\externaldocument[obsolete-]{obsolete}
\externaldocument[fdl-]{fdl}
\externaldocument[index-]{index}

% Theorem environments.
%
\theoremstyle{plain}
\newtheorem{theorem}[subsection]{Theorem}
\newtheorem{proposition}[subsection]{Proposition}
\newtheorem{lemma}[subsection]{Lemma}

\theoremstyle{definition}
\newtheorem{definition}[subsection]{Definition}
\newtheorem{example}[subsection]{Example}
\newtheorem{exercise}[subsection]{Exercise}
\newtheorem{situation}[subsection]{Situation}

\theoremstyle{remark}
\newtheorem{remark}[subsection]{Remark}
\newtheorem{remarks}[subsection]{Remarks}

\numberwithin{equation}{subsection}

% Macros
%
\def\lim{\mathop{\mathrm{lim}}\nolimits}
\def\colim{\mathop{\mathrm{colim}}\nolimits}
\def\Spec{\mathop{\mathrm{Spec}}}
\def\Hom{\mathop{\mathrm{Hom}}\nolimits}
\def\Ext{\mathop{\mathrm{Ext}}\nolimits}
\def\SheafHom{\mathop{\mathcal{H}\!\mathit{om}}\nolimits}
\def\SheafExt{\mathop{\mathcal{E}\!\mathit{xt}}\nolimits}
\def\Sch{\mathit{Sch}}
\def\Mor{\mathop{\mathrm{Mor}}\nolimits}
\def\Ob{\mathop{\mathrm{Ob}}\nolimits}
\def\Sh{\mathop{\mathit{Sh}}\nolimits}
\def\NL{\mathop{N\!L}\nolimits}
\def\CH{\mathop{\mathrm{CH}}\nolimits}
\def\proetale{{pro\text{-}\acute{e}tale}}
\def\etale{{\acute{e}tale}}
\def\QCoh{\mathit{QCoh}}
\def\Ker{\mathop{\mathrm{Ker}}}
\def\Im{\mathop{\mathrm{Im}}}
\def\Coker{\mathop{\mathrm{Coker}}}
\def\Coim{\mathop{\mathrm{Coim}}}

% Boxtimes
%
\DeclareMathSymbol{\boxtimes}{\mathbin}{AMSa}{"02}

%
% Macros for moduli stacks/spaces
%
\def\QCohstack{\mathcal{QC}\!\mathit{oh}}
\def\Cohstack{\mathcal{C}\!\mathit{oh}}
\def\Spacesstack{\mathcal{S}\!\mathit{paces}}
\def\Quotfunctor{\mathrm{Quot}}
\def\Hilbfunctor{\mathrm{Hilb}}
\def\Curvesstack{\mathcal{C}\!\mathit{urves}}
\def\Polarizedstack{\mathcal{P}\!\mathit{olarized}}
\def\Complexesstack{\mathcal{C}\!\mathit{omplexes}}
% \Pic is the operator that assigns to X its picard group, usage \Pic(X)
% \Picardstack_{X/B} denotes the Picard stack of X over B
% \Picardfunctor_{X/B} denotes the Picard functor of X over B
\def\Pic{\mathop{\mathrm{Pic}}\nolimits}
\def\Picardstack{\mathcal{P}\!\mathit{ic}}
\def\Picardfunctor{\mathrm{Pic}}
\def\Deformationcategory{\mathcal{D}\!\mathit{ef}}
\usepackage{cleveref}
\crefname{originalreference}{original reference}{original references}
\Crefname{originalreference}{Original reference}{Original references}
\makeatletter
\newlabel{morphisms-lemma-quasi-affine-finite-type-over-S}{{04II}{source}{}{sourcepointer-1}{}}
\newlabel{morphisms-lemma-quasi-affine-finite-type-over-S@cref}{{[originalreference][0][]04II}{[1][0][]source}{}{}{}}
\makeatother
\crefname{originalreference}{original reference}{original references}
\Crefname{originalreference}{Original reference}{Original references}
\makeatletter
\newlabel{properties-lemma-Noetherian-topology}{{01OZ}{source}{}{sourcepointer-2}{}}
\newlabel{properties-lemma-Noetherian-topology@cref}{{[originalreference][0][]01OZ}{[1][0][]source}{}{}{}}
\makeatother
\crefname{originalreference}{original reference}{original references}
\Crefname{originalreference}{Original reference}{Original references}
\makeatletter
\newlabel{spaces-morphisms-lemma-finite-presentation-noetherian}{{04ZL}{source}{}{sourcepointer-3}{}}
\newlabel{spaces-morphisms-lemma-finite-presentation-noetherian@cref}{{[originalreference][0][]04ZL}{[1][0][]source}{}{}{}}
\makeatother
\crefname{originalreference}{original reference}{original references}
\Crefname{originalreference}{Original reference}{Original references}
\makeatletter
\newlabel{spaces-properties-definition-reduced-induced-space}{{047X}{source}{}{sourcepointer-4}{}}
\newlabel{spaces-properties-definition-reduced-induced-space@cref}{{[originalreference][0][]047X}{[1][0][]source}{}{}{}}
\makeatother
\crefname{originalreference}{original reference}{original references}
\Crefname{originalreference}{Original reference}{Original references}
\makeatletter
\newlabel{spaces-properties-lemma-quasi-separated-sober}{{06NJ}{source}{}{sourcepointer-5}{}}
\newlabel{spaces-properties-lemma-quasi-separated-sober@cref}{{[originalreference][0][]06NJ}{[1][0][]source}{}{}{}}
\makeatother
\crefname{originalreference}{original reference}{original references}
\Crefname{originalreference}{Original reference}{Original references}
\makeatletter
\newlabel{spaces-properties-proposition-locally-quasi-separated-open-dense-scheme}{{06NH}{source}{}{sourcepointer-6}{}}
\newlabel{spaces-properties-proposition-locally-quasi-separated-open-dense-scheme@cref}{{[originalreference][0][]06NH}{[1][0][]source}{}{}{}}
\makeatother
\crefname{originalreference}{original reference}{original references}
\Crefname{originalreference}{Original reference}{Original references}
\makeatletter
\newlabel{topology-lemma-Noetherian}{{0052}{source}{}{sourcepointer-7}{}}
\newlabel{topology-lemma-Noetherian@cref}{{[originalreference][0][]0052}{[1][0][]source}{}{}{}}
\makeatother
\title{Dense open of a finite presented space over a Noetherian base immerses in affine space}
\author{The Stacks Project Authors}
\date{stacks-project:spaces-limits.tex:local-snapshot-20261012; GNU FDL 1.2 or later}
\begin{document}
\maketitle
\noindent\textit{Editorial note.} Counted span is the author's own proof environment, source lines 3313-3366, following the statement at 3305-3311.\par\medskip
\noindent\textit{Editorial note.} The single '(details omitted)' covers the amalgamated immersion A\textasciicircum{}n\_X coprod A\textasciicircum{}m\_X -> A\textasciicircum{}\{max(n,m)+1\}\_X, a standard construction independent of the main argument; flagged in the notes.\par\medskip
\noindent\textit{Editorial note.} Steps the author leaves to the reader are flagged here for the reviewer: source line(s) 3331 carry an "omitted"/"follows easily"/"immediate from the definitions" phrase. Those sentences are retained verbatim; the displayed argument around them stays inside the counted span.\par\medskip
\noindent\textit{Allowed standard prerequisites (editorial).}\begin{itemize}
\item Noetherianity of finite presentation morphisms. Morphisms of Spaces, Lemma spaces-morphisms-lemma-finite-presentation-noetherian (cross-chapter, cited at \textasciitilde{}3314-3316). Y is Noetherian.
\item Dense open subscheme of a quasi-separated space. Properties of Spaces, Proposition spaces-properties-proposition-locally-quasi-separated-open-dense-scheme (cross-chapter, cited at \textasciitilde{}3317-3319 and \textasciitilde{}3359-3360). Reduces Y (resp. X) to a scheme.
\item Noetherian topology decomposition. Topology, Lemma topology-lemma-Noetherian and Properties, Lemma properties-lemma-Noetherian-topology (cross-chapter, cited at \textasciitilde{}3321-3323). Remove intersections of irreducible components.
\item Soberness of quasi-separated spaces. Properties of Spaces, Lemma spaces-properties-lemma-quasi-separated-sober (cross-chapter, cited at \textasciitilde{}3336-3338). T is irreducible with a unique generic point, the image of the generic point of Y.
\item Reduced induced closed subspace. Properties of Spaces, Definition spaces-properties-definition-reduced-induced-space (cross-chapter, cited at \textasciitilde{}3341-3344). The quasi-coherent ideal I cutting out T.
\item Quasi-affine factorization. Morphisms, Lemma morphisms-lemma-quasi-affine-finite-type-over-S (cross-chapter, cited at \textasciitilde{}3362-3364). Factors V -> U as V -> A\textasciicircum{}n\_U -> U.
\end{itemize}\medskip
\ifdefined\hypertarget\hypertarget{sourcepointer-1}{}\fi
\noindent\textit{Source reference.} Original reference 04II: focused-seeds/geometry/S9-spaces-limits-embedding-into-affine-over-qs/ref-tags.txt, source lines 5--5. Navigation pointer only: the counted proof cites this statement by \textbackslash{}ref as a standing input; the deduction the counted proof has to make is carried out inside the counted span, and no part of that deduction is delegated to this pointer.\par\medskip
\ifdefined\hypertarget\hypertarget{sourcepointer-2}{}\fi
\noindent\textit{Source reference.} Original reference 01OZ: focused-seeds/geometry/S9-spaces-limits-embedding-into-affine-over-qs/ref-tags.txt, source lines 6--6. Navigation pointer only: the counted proof cites this statement by \textbackslash{}ref as a standing input; the deduction the counted proof has to make is carried out inside the counted span, and no part of that deduction is delegated to this pointer.\par\medskip
\ifdefined\hypertarget\hypertarget{sourcepointer-3}{}\fi
\noindent\textit{Source reference.} Original reference 04ZL: focused-seeds/geometry/S9-spaces-limits-embedding-into-affine-over-qs/ref-tags.txt, source lines 7--7. Navigation pointer only: the counted proof cites this statement by \textbackslash{}ref as a standing input; the deduction the counted proof has to make is carried out inside the counted span, and no part of that deduction is delegated to this pointer.\par\medskip
\ifdefined\hypertarget\hypertarget{sourcepointer-4}{}\fi
\noindent\textit{Source reference.} Original reference 047X: focused-seeds/geometry/S9-spaces-limits-embedding-into-affine-over-qs/ref-tags.txt, source lines 8--8. Navigation pointer only: the counted proof cites this statement by \textbackslash{}ref as a standing input; the deduction the counted proof has to make is carried out inside the counted span, and no part of that deduction is delegated to this pointer.\par\medskip
\ifdefined\hypertarget\hypertarget{sourcepointer-5}{}\fi
\noindent\textit{Source reference.} Original reference 06NJ: focused-seeds/geometry/S9-spaces-limits-embedding-into-affine-over-qs/ref-tags.txt, source lines 9--9. Navigation pointer only: the counted proof cites this statement by \textbackslash{}ref as a standing input; the deduction the counted proof has to make is carried out inside the counted span, and no part of that deduction is delegated to this pointer.\par\medskip
\ifdefined\hypertarget\hypertarget{sourcepointer-6}{}\fi
\noindent\textit{Source reference.} Original reference 06NH: focused-seeds/geometry/S9-spaces-limits-embedding-into-affine-over-qs/ref-tags.txt, source lines 10--10. Navigation pointer only: the counted proof cites this statement by \textbackslash{}ref as a standing input; the deduction the counted proof has to make is carried out inside the counted span, and no part of that deduction is delegated to this pointer.\par\medskip
\ifdefined\hypertarget\hypertarget{sourcepointer-7}{}\fi
\noindent\textit{Source reference.} Original reference 0052: focused-seeds/geometry/S9-spaces-limits-embedding-into-affine-over-qs/ref-tags.txt, source lines 11--11. Navigation pointer only: the counted proof cites this statement by \textbackslash{}ref as a standing input; the deduction the counted proof has to make is carried out inside the counted span, and no part of that deduction is delegated to this pointer.\par\medskip
\begin{lemma}
\label{lemma-embedding-into-affine-over-qs}
Let $S$ be a scheme. Let $f : Y \to X$ be a morphism of algebraic
spaces over $S$. Assume $X$ Noetherian and $f$ of finite presentation.
Then there exists a dense open $V \subset Y$ and an immersion
$V \to \mathbf{A}^n_X$.
\end{lemma}

\begin{proof}
The assumptions imply that $Y$ is Noetherian
(Morphisms of Spaces, Lemma
\ref{spaces-morphisms-lemma-finite-presentation-noetherian}).
Then $Y$ is quasi-separated, hence has a dense open subscheme
(Properties of Spaces, Proposition
\ref{spaces-properties-proposition-locally-quasi-separated-open-dense-scheme}).
Thus we may assume that $Y$ is a Noetherian scheme.
By removing intersections of irreducible components of $Y$
(use Topology, Lemma \ref{topology-lemma-Noetherian} and
Properties, Lemma \ref{properties-lemma-Noetherian-topology})
we may assume that $Y$ is a disjoint union of irreducible
Noetherian schemes. Since there is an immersion
$$
\mathbf{A}^n_X \amalg \mathbf{A}^m_X
\longrightarrow
\mathbf{A}^{\max(n, m) + 1}_X
$$
(details omitted) we see that it suffices to prove the result in case
$Y$ is irreducible.

\medskip\noindent
Assume $Y$ is an irreducible scheme. Let $T \subset |X|$ be the closure of the
image of $f : Y \to X$. Note that since $|Y|$ and $|X|$ are sober topological
spaces (Properties of Spaces, Lemma
\ref{spaces-properties-lemma-quasi-separated-sober})
$T$ is irreducible with a unique generic point $\xi$ which is the
image of the generic point $\eta$ of $Y$.
Let $\mathcal{I} \subset X$ be a quasi-coherent sheaf of ideals
cutting out the reduced induced space structure on $T$
(Properties of Spaces, Definition
\ref{spaces-properties-definition-reduced-induced-space}).
Since $\mathcal{O}_{Y, \eta}$ is an Artinian local ring we see
that for some $n > 0$ we have $f^{-1}\mathcal{I}^n \mathcal{O}_{Y, \eta} = 0$.
As $f^{-1}\mathcal{I}\mathcal{O}_Y$ is a finite type quasi-coherent ideal
we conclude that $f^{-1}\mathcal{I}^n\mathcal{O}_V = 0$ for
some nonempty open $V \subset Y$. Let $Z \subset X$ be the closed subspace
cut out by $\mathcal{I}^n$. By construction $V \to Y \to X$ factors through
$Z$. Because $\mathbf{A}^n_Z \to \mathbf{A}^n_X$ is an immersion,
we may replace $X$ by $Z$ and $Y$ by $V$.
Hence we reach the situation where $Y$ and $X$ are irreducible and
$Y \to X$ maps the generic point of $Y$ onto the generic point of $X$.

\medskip\noindent
Assume $Y$ and $X$ are irreducible, $Y$ is a scheme,
and $Y \to X$ maps the generic point of
$Y$ onto the generic point of $X$. By Properties of Spaces, Proposition
\ref{spaces-properties-proposition-locally-quasi-separated-open-dense-scheme}
$X$ has a dense open subscheme $U \subset X$. Choose a nonempty affine
open $V \subset Y$ whose image in $X$ is contained in $U$. By
Morphisms, Lemma \ref{morphisms-lemma-quasi-affine-finite-type-over-S}
we may factor $V \to U$ as $V \to \mathbf{A}^n_U \to U$. Composing
with $\mathbf{A}^n_U \to \mathbf{A}^n_X$ we obtain the desired immersion.
\end{proof}

\end{document}
